%% kapa: exam=1 round=32 year=2021 period=2 form=A subject=appraisal_law subject_ko=감정평가관계법규 seq=23 number=23
%% source: Q-Net 감정평가사 기출문제 (cst00309) artlSeq=5208397 file=extracted/1008/5208397/A형/2021년도 제32회 감정평가사 1차 2교시 A형.pdf
%% conversion: kordoc 4.24.51 -> tools/gen1.py (auto); review=proofread-sonnet; flags=none
\begin{question}{23}
국유재산법령상 일반재산에 관한 설명으로 옳은 것은?
\begin{choices}
  \item 정부는 정부출자기업체를 새로 설립하려는 경우에는 일반재산을 현물출자할 수 있다.
  \item 총괄청은 5년 이상 활용되지 아니한 일반재산을 민간사업자와 공동으로 개발할 수 없다.
  \item 중앙관서의 장은 다른 법률에 따라 그 처분이 제한되는 경우에도 일반재산을 매각할 수 있다.
  \item 국가가 직접 행정재산으로 사용하기 위하여 필요한 경우에도 일반재산인 동산과 사유재산인 동산을 교환할 수 없다.
  \item 일반재산을 매각하는 경우 해당 매각재산의 소유권 이전은 매각대금의 완납 이전에도 할 수 있다.
\end{choices}
\end{question}
\answer{1}
